\chapter{LLVM Backend and Linking}
\label{chap:llvm}

The final compiler stage maps lower IR to LLVM IR and asks LLVM to emit a native
object file. This chapter explains the LLVM backend implementation, including
function declaration, type mapping, value mapping, basic block construction, phi
node construction, object emission, and linking.

\section{Why LLVM}

Writing a complete native backend is a large project. A backend must understand
instruction selection, register allocation, stack layout, object file formats,
relocations, calling conventions, debug information, and target-specific
details. LLVM provides reusable infrastructure for many of these tasks
\cite{lattner2004llvm,llvm_langref}. By targeting LLVM IR, \epl{} can focus on
source-language design and compiler lowering while still producing native object
files.

The project originally considered multiple possible backends in its proposal,
including LLVM and self-contained minimal backends. The implemented compiler
uses LLVM because it gives the project an end-to-end native-code path within the
time available for a final year project.

\section{LLVM Module Construction}

The LLVM backend is implemented in \texttt{src/llvm.rs}. The main public
function is \texttt{LlvmModule::from\_ir}. It takes lower IR and returns an LLVM
module wrapper. The wrapper owns an LLVM context and LLVM module pointer and
disposes of them when dropped.

LLVM construction happens in two phases. First, the compiler declares every
function and records its LLVM function type and LLVM function value. This allows
calls to functions defined later in the file. Second, the compiler emits the body
of each function that has a body. External declarations are left as LLVM
declarations.

This two-phase structure mirrors semantic analysis. Function signatures are
known before bodies are emitted. This is necessary for recursion and forward
calls.

\section{Type Mapping}

The LLVM type mapping is direct and is shown in Table~\ref{tab:llvm-types}.

\begin{table}[H]
  \centering
  \begin{tabular}{p{0.30\textwidth}p{0.45\textwidth}}
    \toprule
    Lower IR type & LLVM type \\
    \midrule
    \texttt{Unit} & empty LLVM struct \\
    \texttt{Bool} & \texttt{i1} \\
    \texttt{I8} & \texttt{i8} \\
    \texttt{I32} & \texttt{i32} \\
    \texttt{I64} & \texttt{i64} \\
    \texttt{Ptr} & opaque LLVM pointer \\
    \texttt{Struct} & LLVM struct with lowered field types \\
    \texttt{Array} & LLVM array with lowered element type and length \\
    \bottomrule
  \end{tabular}
  \caption{Mapping from lower IR types to LLVM types}
  \label{tab:llvm-types}
\end{table}

LLVM uses a distinct \texttt{i1} type for booleans. Integer widths map
directly. Pointers use LLVM opaque pointers. Unit values become an empty struct,
although lower IR cleanup removes most zero-sized operations before LLVM
generation.

\section{Values}

The backend maintains a map from lower IR definitions to LLVM values. Function
arguments are inserted into the map from the LLVM function parameters. Allocation
slots are emitted in the entry block using LLVM \texttt{alloca} and inserted
into the map. Instruction results are inserted as instructions are emitted.
Block arguments are inserted as LLVM phi nodes before block instructions are
emitted.

Constants are translated as needed. Boolean constants become LLVM integer
constants of type \texttt{i1}. Number constants become LLVM integer constants of
the correct width. Undefined values become LLVM \texttt{undef}. String constants
are emitted as global arrays initialized with the string data and null
terminator.

\section{Basic Blocks and Phi Nodes}

For each lower IR basic block, the backend creates an LLVM basic block. It then
creates phi nodes for non-entry block arguments. The actual incoming values are
filled when predecessor terminators are emitted.

This maps lower IR block arguments to LLVM's phi-node representation. For a
jump terminator, the backend emits an LLVM branch and then adds incoming values
to the target block's phi nodes. For a conditional jump, it emits an LLVM
conditional branch and adds incoming values to the phi nodes of both successors.

This design keeps lower IR independent from LLVM's exact phi instruction API
while still preserving SSA-style value merging.

\section{Instruction Emission}

Each lower IR instruction maps to one or more LLVM builder calls:

\begin{itemize}
  \item \texttt{Load} maps to \texttt{LLVMBuildLoad2};
  \item \texttt{Store} maps to \texttt{LLVMBuildStore};
  \item \texttt{FunctionCall} maps to \texttt{LLVMBuildCall2};
  \item \texttt{Cmp} maps to \texttt{LLVMBuildICmp} with signed or unsigned
        predicates;
  \item \texttt{Arithmetic} maps to add, sub, mul, signed/unsigned division, or
        signed/unsigned remainder;
  \item \texttt{Not} maps to XOR with \texttt{true};
  \item \texttt{OffsetPtr} maps to \texttt{LLVMBuildGEP2} over \texttt{i8};
  \item \texttt{Zext}, \texttt{Sext}, and \texttt{Truncate} map to LLVM integer
        cast instructions.
\end{itemize}

Signedness is handled at the instruction level. For example, a lower IR
comparison stores whether the operands should be interpreted as signed. The LLVM
backend chooses \texttt{LLVMIntSLT} or \texttt{LLVMIntULT} for less-than based
on that flag.

\section{Terminator Emission}

Lower IR terminators map directly to LLVM terminators. A jump becomes an LLVM
branch. A conditional jump becomes an LLVM conditional branch. A return becomes
an LLVM return with a value. An unreachable terminator becomes LLVM
\texttt{unreachable}.

The backend must ensure that each LLVM basic block ends with exactly one
terminator. This is one reason lower IR separates instructions from terminators.
The structure prevents accidental emission of non-terminator instructions after
a branch or return.

\section{Module Verification}

After constructing the LLVM module, the compiler calls \texttt{LLVMVerifyModule}.
If verification fails, the compiler prints the LLVM verifier error and exits.
Verification is important because LLVM IR has structural rules that are easy to
violate when building IR through an API. Examples include missing terminators,
phi nodes with incorrect incoming blocks, mismatched call signatures, or invalid
types.

LLVM verification does not prove that the original \epl{} program is safe or
that the generated code is semantically correct. It only checks LLVM IR
well-formedness. It is nevertheless an important backend safety net.

\section{Object File Emission}

The \texttt{llvm-obj} command calls \texttt{LlvmModule::compile}. The compile
method initializes LLVM targets, obtains the default target triple, creates a
target machine, sets the module data layout and target triple, and asks LLVM to
emit an object file named \texttt{a.out.o}.

The compiler currently emits object files with no LLVM optimization pipeline
enabled. There is commented code showing where a pass pipeline could be run in
the future. This is a reasonable current decision because the project focuses on
correct lowering rather than performance optimization.

\section{Linking}

The compiler stops at an object file. To produce an executable, the object file
is linked with a system C compiler:

\begin{lstlisting}[caption={Compiling and linking an EPL program}]
cargo run -- llvm-obj examples/hello_world.epl
cc a.out.o
./a.out
\end{lstlisting}

This link step resolves external declarations such as \texttt{printf} and
\texttt{exit}. It also lets the platform C toolchain handle startup files and
libc linkage. A future compiler could invoke the linker automatically, but
keeping the link step explicit makes the current architecture simpler.

\section{FFI Limits}

Although \epl{} can call simple external C functions, it should not be described
as having a complete C FFI. The implemented path works for examples such as
\texttt{printf(fmt: *i8, ...)} and \texttt{exit(code: i32) -> !}. It does not
fully model all C ABI cases.

Important missing cases include by-value aggregate arguments and returns,
platform-specific varargs details, exact C \texttt{void} mapping for every
signature shape, structure packing pragmas, integer promotion rules for
varargs, callbacks, and library header parsing. These features are not needed to
demonstrate the compiler pipeline, but they are important for future work if
\epl{} is to call arbitrary C libraries reliably.

\section{Backend Tradeoffs}

The LLVM backend is direct and readable. It does not attempt to hide LLVM behind
a large abstraction layer. The advantage is that every lower IR construct has an
obvious LLVM emission point. The disadvantage is that target-specific details
are partly delegated to LLVM without building a higher-level ABI model inside
\epl{}.

For a final year compiler project, this is an appropriate tradeoff. It produces
native object code, gives access to LLVM verification, and keeps the backend
small enough to reason about.
